%%
%% Full version (appendices) of the PVLDB Volume 20 paper; a technical report in the artifact repository.
%% Same results bridge as the paper: every measured number is \res{key}.
%%
\documentclass[sigconf,nonacm,review,balance=false]{acmart}

\settopmatter{printacmref=false,printfolios=true}

%% Shared preamble of the main paper and the appendix.
\usepackage{algorithm}
\usepackage{algpseudocode}
\usepackage{multirow}
\usepackage{graphicx}
\usepackage{enumitem}
\usepackage{needspace}
\usepackage{colortbl}
%% Table highlighting: ours (blue, the colour of c = 1.1 in the figures), reference (grey), failure or cost (red).
\definecolor{hlours}{HTML}{DCEBF7}
\definecolor{hlref}{HTML}{ECECEC}
\definecolor{hlbad}{HTML}{F9DEDC}
\newcommand{\hb}{\cellcolor{hlbad}}
%% Lists: labels flush with the paragraph left edge, no indentation.
\setlist[itemize,enumerate]{leftmargin=*,labelindent=0pt,itemindent=0pt,listparindent=0pt,topsep=2pt,partopsep=0pt}
\setlist[description]{leftmargin=0pt,labelindent=0pt,itemindent=0pt}
\graphicspath{{figures/}}

\newcommand{\sys}{\textsc{Ariadne}}
\newcommand{\Adm}{\mathrm{Adm}}
\newcommand{\cser}{$c$-serializable}
\newcommand{\dist}{d}
\newcommand{\Ds}{D_{\mathrm{s}}}
\newcommand{\rhoa}{\rho}
\newcommand{\Rc}{R_c}
\newcommand{\Read}{\textsc{Read}}
\newcommand{\pol}[1]{\textsf{#1}}
\algrenewcommand\algorithmicindent{0.9em}
\emergencystretch=1.5em

%% ---------------------------------------------------------------------
%% Results bridge.  Every measured number in the text is \res[shape]{key}
%% (or \ra, \rb, \ri, \rarng).  paper/make_results.py computes the keys
%% from the result files and writes results.tex (\resdef{key}{value}).
%% A key without a value prints its shape in orange in a draft; with
%% \resultsfinaltrue it is a compile error, so no missing number can
%% survive into a final build.
%% ---------------------------------------------------------------------
\newif\ifresultsfinal
\resultsfinalfalse
\newwrite\reskeys
\immediate\openout\reskeys=\jobname.reskeys
\newcommand{\resdef}[2]{\expandafter\gdef\csname res/#1\endcsname{#2}}
\DeclareRobustCommand{\res}[2][XX]{%
  \immediate\write\reskeys{#2}%
  \ifcsname res/#2\endcsname\csname res/#2\endcsname\else
    \ifresultsfinal\PackageError{ariadne}{No result for key #2}{Run paper/make_results.py on the results.}\fi
    {\color{orange}#1}%
  \fi}
%% A value computed from smoke-scale results (make_results.py --suffix _smoke): orange in a draft, an error in a final build.
\newcommand{\resprov}[1]{\ifresultsfinal\PackageError{ariadne}{Smoke-scale value in a final build}{Run paper/make_results.py on the full results.}\fi{\color{orange}#1}}
\newcommand{\ra}[1]{\res[00.0]{#1}}   % one decimal
\newcommand{\rb}[1]{\res[0.00]{#1}}   % two decimals
\newcommand{\ri}[1]{\res[00]{#1}}     % integer
%% A range of two keys at one decimal; prints one value when both round to the same printed number.
\newcommand{\rarng}[2]{\immediate\write\reskeys{#2}%
  \ifcsname res/#1\endcsname\ifcsname res/#2\endcsname
    \expandafter\let\expandafter\rarngA\csname res/#1\endcsname\expandafter\let\expandafter\rarngB\csname res/#2\endcsname
    \ifx\rarngA\rarngB\ra{#1}\else\ra{#1} to \ra{#2}\fi
  \else\ra{#1} to \ra{#2}\fi\else\ra{#1} to \ra{#2}\fi}
%% Figures are drawn from the results by ariadne.plots (paper/build.sh).  A figure that is not there yet is a framed
%% note in a draft and an error in a final build.
\newcommand{\figstatus}[1]{}
\newcommand{\datafig}[2][\columnwidth]{\IfFileExists{figures/#2.pdf}{\includegraphics[width=#1]{#2}}%
  {\ifresultsfinal\PackageError{ariadne}{Missing figure #2}{Run paper/build.sh on the results.}\fi
   \fbox{\parbox[c][0.45#1][c]{\dimexpr#1-2\fboxsep-2\fboxrule}{\centering\color{orange}\sffamily\small figure \texttt{\detokenize{#2}}: pending}}}}
\InputIfFileExists{results.tex}{}{}

%% The appendix is the full version of the paper (a separate technical report, in the artifact repository).
%% In the 12-page paper the appendix sections are cited as sections of the full version.
\newcommand{\appProofs}{the full version~\cite[App.~A]{ariadnefull}}
\newcommand{\appPG}{the full version~\cite[App.~B]{ariadnefull}}
\newcommand{\appLaw}{the full version~\cite[App.~C]{ariadnefull}}
\newcommand{\appExp}{the full version~\cite[App.~D]{ariadnefull}}
\newcommand{\appAudit}{the full version~\cite[App.~E]{ariadnefull}}
\newcommand{\appRepro}{the full version~\cite[App.~F]{ariadnefull}}
\newcommand{\appRelated}{the full version~\cite[App.~G]{ariadnefull}}
\newcommand{\AppProofs}{The full version~\cite[App.~A]{ariadnefull}}
\newcommand{\AppExp}{The full version~\cite[App.~D]{ariadnefull}}

\usepackage{xr}\externaldocument[M-]{pvldb20-ariadne}% table and figure numbers of the 12-page paper
\usepackage{adjustbox}
\usepackage{fancyvrb}
\usepackage{placeins}
\usepackage{stfloats}
%% The appendix is mostly figures: allow dense float placement and float pages.
\renewcommand{\topfraction}{0.95}\renewcommand{\dbltopfraction}{0.95}
\renewcommand{\bottomfraction}{0.6}\renewcommand{\textfraction}{0.05}
\renewcommand{\floatpagefraction}{0.50}\renewcommand{\dblfloatpagefraction}{0.50}
\setcounter{topnumber}{6}\setcounter{dbltopnumber}{4}\setcounter{bottomnumber}{3}\setcounter{totalnumber}{9}
%\resultsfinaltrue

\newlength{\afw}
%% An auto-generated table (paper/tables/*.tex, written by ariadne.plots.tables) at its natural size:
%% single column if it fits, two columns otherwise, scaled down only if it is wider than the page.
%% \autotab{name}{label}{caption}
\newcommand{\tabstatus}[1]{}
\newcommand{\autotab}[3]{%
  \IfFileExists{tables/#1.tex}{\settowidth{\afw}{\small\input{tables/#1.tex}}}{\setlength{\afw}{0pt}}%
  \ifdim\afw=0pt
    \ifresultsfinal\PackageError{ariadne}{Missing table #1}{Run paper/build.sh}\fi
    \begin{table}[!tbp]\caption{#3}\label{#2}\centering{\color{orange}[table \texttt{\detokenize{#1}}: pending]}\end{table}%
  \else\ifdim\afw>1.02\columnwidth
    \begin{table*}[!tbp]\caption{#3}\label{#2}\centering\small
      \begin{adjustbox}{max width=\textwidth, max totalheight=0.82\textheight}\input{tables/#1.tex}\end{adjustbox}\tabstatus{#1}\end{table*}%
    \tabcont{#1}{#2}{2}\tabcont{#1}{#2}{3}\tabcont{#1}{#2}{4}\tabcont{#1}{#2}{5}\tabcont{#1}{#2}{6}%
  \else
    \begin{table}[!tbp]\caption{#3}\label{#2}\centering\small\begin{adjustbox}{max width=\columnwidth, max totalheight=0.82\textheight}\input{tables/#1.tex}\end{adjustbox}\tabstatus{#1}\end{table}%
  \fi\fi}

%% continuation of a generated table split by tidy_tables.py (tables/<name>__<k>.tex)
\newcommand{\tabcont}[3]{\IfFileExists{tables/#1__#3.tex}{%
  \begin{table*}[!tbp]\centering\small{\sffamily\bfseries Table~\ref{#2} (continued)}\par\smallskip
  \begin{adjustbox}{max width=\textwidth, max totalheight=0.88\textheight}\input{tables/#1__#3.tex}\end{adjustbox}\end{table*}}{}}

%% A data figure at its natural size: single column if it fits (never scaled up), two columns otherwise.
%% \autofig{file}{label}{description}{caption}
\newcommand{\autofig}[4]{%
  \IfFileExists{figures/#1.pdf}{\settowidth{\afw}{\includegraphics{#1}}}{\setlength{\afw}{0pt}}%
  \ifdim\afw=0pt
    \ifresultsfinal\PackageError{ariadne}{Missing figure #1}{Run paper/build.sh on the final results.}\fi
    \begin{figure}[!tbp]\centering{\color{orange}[figure \texttt{\detokenize{#1}}: pending, needs the full results]}\caption{#4}\label{#2}\end{figure}%
  \else\ifdim\afw>1.3\columnwidth
    \ifdim\afw>\textwidth\setlength{\afw}{\textwidth}\fi
    \begin{figure*}[!tbp]\centering\begin{adjustbox}{max totalheight=0.78\textheight}\datafig[\afw]{#1}\end{adjustbox}\Description{#3}\caption{#4}\label{#2}\end{figure*}%
  \else
    \ifdim\afw>\columnwidth\setlength{\afw}{\columnwidth}\fi
    \begin{figure}[!tbp]\centering\datafig[\afw]{#1}\Description{#3}\caption{#4}\label{#2}\end{figure}%
  \fi\fi}
%% the same with top placement only, for a figure that must share its page with text
\newcommand{\autofigtop}[4]{%
  \IfFileExists{figures/#1.pdf}{\settowidth{\afw}{\includegraphics{#1}}}{\setlength{\afw}{0pt}}%
  \ifdim\afw=0pt
    \ifresultsfinal\PackageError{ariadne}{Missing figure #1}{Run paper/build.sh on the final results.}\fi
    \begin{figure}[!t]\centering{\color{orange}[figure \texttt{\detokenize{#1}}: pending, needs the full results]}\caption{#4}\label{#2}\end{figure}%
  \else\ifdim\afw>1.3\columnwidth
    \ifdim\afw>\textwidth\setlength{\afw}{\textwidth}\fi
    \begin{figure*}[!t]\centering\begin{adjustbox}{max totalheight=0.78\textheight}\datafig[\afw]{#1}\end{adjustbox}\Description{#3}\caption{#4}\label{#2}\end{figure*}%
  \else
    \ifdim\afw>\columnwidth\setlength{\afw}{\columnwidth}\fi
    \begin{figure}[!t]\centering\datafig[\afw]{#1}\Description{#3}\caption{#4}\label{#2}\end{figure}%
  \fi\fi}

\begin{document}

\title{Admissible, Not Repeatable: Isolating Approximate Similarity Reads in Shared Agent Memory (Full Version: Appendices)}

\author{Chang Lu}
\affiliation{%
  \department{Computational Biology \& Biomedical Informatics Program}
  \institution{Yale University}
  \city{New Haven, CT}
  \country{USA}
}
\email{chang.lu@yale.edu}

\author{Yizheng Zhao}
\affiliation{%
  \department{State Key Laboratory of Novel Software Technology}
  \institution{Nanjing University}
  \city{Nanjing}
  \country{China}
}
\email{zhaoyz@nju.edu.cn}

\maketitle

\noindent
This technical report holds the appendices of the paper of the same
title, submitted to PVLDB Volume~20. Part~I (\S\ref{app:proofs}) contains
the proofs and the precise statements behind \S\S4 and~5 of the
paper. Part~II (\S\S\ref{app:pg} to~\ref{app:repro}) documents the
experiments that the paper summarizes in one sentence each. Appendix~\ref{app:related} discusses related work, and Appendix~\ref{app:notation} collects all symbols and terms. Section,
theorem, figure, and table numbers without a letter are the paper's.

\appendix
\renewcommand{\thefigure}{A\arabic{figure}}
\renewcommand{\thetable}{A\arabic{table}}

%% ====================================================================
\section{Proofs and Precise Statements}
\label{app:proofs}

\subsection{Conventions}
\label{app:conv}

States are finite sets of items whose vectors lie in a metric space
with distance $\dist$; $\delta_i(S,q)$ is the $i$-th smallest distance
from $q$ to $S$, and reads ask for $k\le|S|$ items, as in
\S2. For the exact
statements below we assume that the distances from a query to the items
of a state are pairwise distinct; with ties, the statements hold for
any fixed tie-breaking rule, with ``$<$'' between tied distances read
as ``precedes.'' A transaction is a deterministic function from the
answers it has received so far to its next read or to its write set;
Algorithm~1 fixes the queries of a transaction before its first read,
and the theorems below are stated for such transactions. A
write set consists of a set of insertions, each of a new item, and a
set of deletions of items in the transaction's held answers (a
replacement contributes to both). It is applied atomically with respect
to reads: a read observes either all of it or none of it. On stores
without multi-statement transactions this holds for write sets with one
operation; for larger write sets \sys{} relies on the commit lock and on
reads that are validated under it. A read-only transaction with one
read is not validated, so on such stores Theorem~\ref{thm:safety-full}
needs write sets of one operation, or validation of those transactions
as well.

\begin{lemma}[Final state]
\label{lem:final}
Let $I$ and $U$ be the items inserted and deleted by the transactions
of an execution. Every order $\pi$ that satisfies Definition~4.1 for
some $c$ has last state $(S_0\cup I)\setminus U$.
\end{lemma}
\begin{proof}
An item $y\in U$ deleted by the $j$-th transaction in $\pi$ lies in a
held answer of that transaction, hence in $S^\pi_{j-1}$, because
admissible answers consist of items of the state. So every deletion
removes an item that is present, and no item is deleted twice: after
its deletion it is absent from every later state, and a later held
answer containing it would not be admissible. Every insertion adds a
new item, which no earlier transaction can delete. Hence each item of
$S_0\cup I$ is present at the end exactly if it is not in $U$.
\end{proof}

\subsection{Operational reading of Definition~4.1}
\label{app:oper}

Call a store a \emph{$c$-store} if its similarity read $\Read(q,k)$ in
state $S$ may return any element of $\Adm_c(S,q,k)$.

\begin{proposition}
\label{prop:oper}
An execution is \cser{} if and only if there is a total order $\pi$ of
its transactions such that running the transactions one at a time in
the order $\pi$ on a $c$-store, with every read returning its held
answer in the execution, is a legal run in which every transaction
produces its write set in the execution.
\end{proposition}
\begin{proof}
Suppose the execution is \cser{} with order $\pi$. We show by
induction on $j$ that the serial run reaches state $S^\pi_{j}$ after
its $j$-th transaction. The $j$-th transaction starts in
$S^\pi_{j-1}$. Its writes are buffered, so all its reads are evaluated
in $S^\pi_{j-1}$, and by Definition~4.1 each held answer is in
$\Adm_c(S^\pi_{j-1},q,k)$, which a $c$-store may return. Shown the same
answers, the transaction issues the same reads and the same write set
by determinism, and the run reaches $S^\pi_j$. Conversely, a legal
serial run with the held answers shows
$A\in\Adm_c(S^\pi_{j-1},q,k)$ for every read directly.
\end{proof}

For $c=1$ and distinct distances $\Adm_1(S,q,k)$ is a singleton, a
$1$-store is an exact store, and the proposition is view
serializability restricted to transactions whose reads are similarity reads.

\subsection{Separation (Proposition~4.2) is tight}
\label{app:sep}

\begin{proposition}
\label{prop:tight}
For every $\tau>0$, $c>1$ and $\varepsilon>0$ there is, on the real
line with $\dist(x,y)=|x-y|$, a \cser{} execution of one
$\tau$-deduplicating insertion on a $(\tau/c)$-separated state whose
final state contains two items at distance less than
$\tau/c+\varepsilon$.
\end{proposition}
\begin{proof}
Let $k=1$, $0<\eta<\varepsilon$, the query $q=0$, and
$S_0=\{a,b\}$ with $a=\tau/c+\eta$ and $b=-(\tau+\eta)$. The two items
are farther apart than $\tau/c$. The exact answer of $\Read(q,1)$ is $(a)$,
and the list $(b)$ is $c$-admissible because
$\dist(q,b)=\tau+\eta\le \tau+c\,\eta=c\,\delta_1(S_0,q)$. A
transaction that observes $(b)$ sees $\dist(q,a_1)>\tau$ and inserts
$q$; the final state contains $a$ and $q$ at distance
$\tau/c+\eta<\tau/c+\varepsilon$.
\end{proof}

With a validation read of approximation ratio $\rho$,
Theorem~4.3 gives $(c\rho)$-serializability and therefore, if $S_0$
is $(\tau/(c\rho))$-separated, $(\tau/(c\rho))$-separation; the certified variant ($\rho=1$) gives
$\tau/c$. \S\ref{app:exp-correct} reports the smallest stored
pairwise distances that we measured in the certified runs against these bounds.

\subsection{Safety (Theorem~4.3) for general transactions}
\label{app:safety}

We restate the theorem for transactions with several reads and with
insertions, deletions, and replacements, and construct the serial order
explicitly from the order of the commit lock.

\begin{theorem}
\label{thm:safety-full}
Consider an execution in which every transaction runs Algorithm~1 under
one commit lock, with validation reads issued after the lock is
granted, on a store whose reads observe every write set applied before
they are issued and none in part, and suppose every answer that the store returns is
$\rho$-admissible in the state it is evaluated in. Then the execution
is $(c\rho)$-serializable. A transaction that aborts at the round limit
writes nothing and is not part of the execution.
\end{theorem}
\begin{proof}
Number the transactions that take the commit lock as $T_1,T_2,\dots$ in
the order in which they hold it in their last iteration,
and let $S_j$ be the state after $T_j$ applied its write set, with
$S_0$ the initial state, before any write set is applied.

\emph{The validation state.} All write sets are applied while the
commit lock is held, and the lock is held from the validation reads to
the application of the writes. Hence between the moment $T_j$ acquires
the lock in its last iteration and the moment it applies its writes no
other write set is applied, and all write sets applied before are those
of $T_1,\dots,T_{j-1}$. The validation reads of $T_j$ are therefore
evaluated in $S_{j-1}$, and $T_j$'s write set transforms $S_{j-1}$ into
$S_j$.

\emph{Admissibility.} Let $A$ be the answer of a read $\Read(q,k)$ that
$T_j$ used in its last call of \textsc{Think}, and $F$ the
corresponding fresh answer. The test $A\sqsubseteq_c F$ gives
$A\subseteq S_{j-1}$, $|A|=|F|=k$, and
$\dist(q,a_i)\le c\,\dist(q,f_i)$ for all $i$. Since
$\rho_F=\rhoa(F\mid S_{j-1},q)\le\rho$,
\[
\dist(q,a_i)\;\le\; c\,\dist(q,f_i)\;\le\; c\,\rho_F\,\delta_i(S_{j-1},q),
\]
so $A\in\Adm_{c\rho}(S_{j-1},q,k)$ and
$\rhoa(A\mid S_{j-1},q)\le c\,\rho_F$. The write set of $T_j$ is the
one that \textsc{Think} computed from these answers. The bound of
Theorem~4.3 is attained: with $k=1$, a nearest item at distance $1$
from $q$, a fresh answer $F=(f_1)$ at distance $\rho_F$, and a
held answer $A=(a_1)$ at distance $c\rho_F$, the held answer passes
the test with equality at $i=1$, and its ratio is $c\,\rho_F$.

\emph{Read-only transactions with one read.} Such a transaction $T'$
observes an answer that the store evaluated in some state. Write sets
are applied atomically with respect to reads (\S\ref{app:conv}), so this state is $S_j$
for some $j$, and the answer is $\rho$-admissible in it. Insert $T'$
into the order directly after $T_j$. It changes no state, so the
states that the order assigns to the other transactions are unchanged.

The resulting total order explains every read by a state in which its
answer is $(c\rho)$-admissible, as Definition~4.1 requires.
\end{proof}

\paragraph{Remarks}
(1)~Earlier iterations of the loop, and the answers they used, do not
appear in the execution: only the write set of the last iteration is
applied. (2)~The proof uses the commit lock twice: all writers must take
it, and the validation reads must observe the state at the time the
lock is held. A lock per group of agents on a shared store violates the
first requirement. A validation transaction whose snapshot was taken
before the lock was granted violates the second; this is what happens
on PostgreSQL at repeatable read, where the statement that requests the
advisory lock fixes the snapshot (\S\ref{app:audit}). (3)~The
theorem bounds the ratio of the held answer by $c\,\rho_F$ for
the \emph{realized} ratio $\rho_F$ of each validation read, which is
what the auditor checks commit by commit. (4)~A read-only transaction
with several reads is not serialized by its reads, because writes may
commit between them; it is validated like any other transaction.

\subsection{Characterization (Theorem~5.1) and the conflict radius}
\label{app:char}

\begin{lemma}[General form]
\label{lem:general}
Let $A$ be the exact answer of $\Read(q,k)$ in $S$, with distances
$r_1<\dots<r_k$, and let $S'=(S\cup V)\setminus U$ for a set $V$ of
inserted items and a set $U$ of deleted items. Then
$A\in\Adm_c(S',q,k)$ if and only if $A\cap U=\emptyset$ and
$r_i\le c\,\delta_i(S',q)$ for all $i\le k$. If $A\cap U=\emptyset$,
then $A\in\Adm_c(S\cup V,q,k)$ implies $A\in\Adm_c(S',q,k)$; in
particular, deletions alone never make $A$ inadmissible unless they
delete an item of $A$.
\end{lemma}
\begin{proof}
The first claim is Definition~2.1. For the second, removing items from
a state cannot decrease any $\delta_i$, so
$\delta_i(S',q)\ge\delta_i(S\cup V,q)$. For $V=\emptyset$ and
$A\cap U=\emptyset$, the items of $A$ are the $k$ nearest items of $S$
and belong to $S\setminus U$, so they are its $k$ nearest items as
well, and $A$ is the exact answer in $S'$.
\end{proof}

Theorem~5.1 is the case $V=\{v\}$, $U=\emptyset$: with
$t=\dist(q,v)<r_k$ and rank $p=1+|\{i:r_i<t\}|$, the sorted distances
of $S\cup\{v\}$ are $r_1,\dots,r_{p-1},t,r_p,\dots$, so
$\delta_i=r_i$ for $i<p$, $\delta_p=t$, and $\delta_i=r_{i-1}$ for
$p<i\le k$; Lemma~\ref{lem:general} then gives Theorem~5.1.

\begin{proof}[Proof of Corollary~5.2]
Put $r_0=0$ and let $J=\{i\in\{1,\dots,k\}: r_i>c\,r_{i-1}\}$. The index
$1$ belongs to $J$, so $j=\max J$ exists and $\Rc(q)=r_j/c$, because
the $r_i$ increase. Let a new item $v$ be inserted at distance $t=\dist(q,v)$ from the query $q$.

If $t\ge r_k$, the top-$k$ is unchanged, $A$ stays admissible, and
$t\ge r_k\ge r_j/c=\Rc(q)$, because $j\le k$ and $c\ge1$.

Otherwise let $p\le k$ be the rank of $v$, so that
$r_{p-1}\le t\le r_p$. If $p<j$, the second condition of Theorem~5.1
fails at $i=j$, so $A$ is inadmissible, and
$t\le r_p\le r_{j-1}<r_j/c=\Rc(q)$. If $p=j$, no index $i>p$ lies in
$J$, so the second condition holds and $A$ is inadmissible exactly if
the first fails, that is, if $t<r_j/c=\Rc(q)$. If $p>j$, again the
second condition holds, and $p\notin J$ gives
$r_p\le c\,r_{p-1}\le c\,t$, so the first holds too and $A$ is
admissible; moreover $t\ge r_{p-1}\ge r_j\ge\Rc(q)$. In every case $A$
is inadmissible if and only if $t<\Rc(q)$. The bounds
$r_1/c\le\Rc(q)\le r_k/c$ follow from $1\in J$ and $j\le k$.
\end{proof}

\paragraph{Properties}
$\Rc(q)$ is non-increasing in $c$, and $R_1(q)=r_k$ for distinct
distances: under strict semantics the whole top-$k$ ball is protected,
which is the predicate that re-execution validation enforces on an
exact index. If all consecutive ratios $r_i/r_{i-1}$, $i\ge2$, are at
most $c$, then $\Rc(q)=r_1/c$: only an insertion that becomes the new
nearest neighbor by a factor of more than $c$ conflicts.

\paragraph{Several insertions}
The radius form does not extend to $|V|>1$. For $k=3$,
$r=(1,1.09,1.18)$ and $c=1.1$ we have $\Rc(q)=1/1.1\approx0.909$.
Two insertions at distances $0.95$ and $0.96$ both lie outside
$\Rc(q)$, yet the new second-smallest distance is $0.96$ and
$r_2=1.09>1.1\cdot0.96$. The general test of Lemma~\ref{lem:general}
applies, and the conflict law treats several insertions by simulation
(\S\ref{app:law-proof}).

\subsection{Conflict law (Theorem~5.3)}
\label{app:law-proof}

\paragraph{Model}
Around the query $q$ the stored items form a Poisson process whose
expected number of points within distance $s$ is $\lambda s^{D}$ on
the range of distances that matters for the read (the model is used for
the $k+1$ nearest items). In the \emph{mass coordinate}
$\omega=\lambda s^{D}$ the process is homogeneous with unit rate, and a
distance ratio $c$ corresponds to the mass ratio $c^{D}=1/x$. Let
$\mu_1<\mu_2<\dots$ be the masses of the stored neighbors and
$\mu_0=0$. A concurrent in-ball insertion with dimension ratio
$\alpha=\Ds/D\in(0,1]$ has mass $\omega$ with density proportional to
$\omega^{\alpha-1}$ on $(0,\mu_k)$, independently of the ratios $B_i$
of consecutive neighbor masses defined next.

\begin{lemma}
\label{lem:beta}
The ratios $B_i=\mu_i/\mu_{i+1}$, $i\ge1$, are independent, $B_i$
has density $i\,b^{i-1}$ on $[0,1]$, and $(B_1,\dots,B_{k-1})$ is
independent of $\mu_k$.
\end{lemma}
\begin{proof}
Given $\mu_{i+1}$, the points $\mu_1,\dots,\mu_i$ are the
order statistics of $i$ independent uniform variables on
$(0,\mu_{i+1})$, so $\mu_i/\mu_{i+1}$ is the maximum of $i$ uniform
variables on $(0,1)$, with density $i\,b^{i-1}$, and
$(\mu_1,\dots,\mu_i)/\mu_{i+1}$ is independent of $\mu_{i+1}$ and of
the process beyond it. Applying this for
$i=k-1,k-2,\dots,1$ gives the claim.
\end{proof}

\begin{proof}[Proof of Theorem~5.3]
Because the density of $\omega$ on $(0,\mu_k)$ is
$\alpha \omega^{\alpha-1}/\mu_k^{\alpha}$, for an event $E$ determined by
$\omega$ and the $\mu_i$ we have
$\Pr[E]=\mathbb{E}\bigl[\mu_k^{-\alpha}\int_E\alpha
\omega^{\alpha-1}\,d\omega\bigr]$. In mass coordinates Theorem~5.1 reads: the
insertion has rank $p$ and the answer stays $c$-admissible if and only
if
\[
\max(\mu_{p-1},\,x\,\mu_p)\le \omega<\mu_p
\quad\text{and}\quad B_i\ge x\ \text{ for } p\le i<k .
\]
With $B_0=0$, integrating $\alpha \omega^{\alpha-1}$ over the interval for
$\omega$ gives
\[
\mu_p^{\alpha}\bigl(1-\max(B_{p-1},x)^{\alpha}\bigr),
\qquad \mu_p=\mu_k\textstyle\prod_{i=p}^{k-1}B_i .
\]
With the factor $\mu_k^{-\alpha}$ and the conditions $B_i\ge x$ for $p\le i<k$, the expectation is
\begin{align*}
\varphi_\alpha(p)
&=\mathbb{E}\Bigl[\bigl(1-\max(B_{p-1},x)^{\alpha}\bigr)
   \prod_{i=p}^{k-1}B_i^{\alpha}\,\mathbf{1}\{B_i\ge x\}\Bigr]\\
&=\mathbb{E}\bigl[1-\max(B_{p-1},x)^{\alpha}\bigr]\;
  \prod_{i=p}^{k-1}\mathbb{E}\bigl[B_i^{\alpha}\mathbf{1}\{B_i\ge x\}\bigr]
\end{align*}
by Lemma~\ref{lem:beta}. For $i\ge1$,
$\mathbb{E}[B_i^{\alpha}\mathbf{1}\{B_i\ge x\}]
=\int_x^1 i\,b^{i-1+\alpha}\,db=\frac{i}{i+\alpha}(1-x^{i+\alpha})$.
For $p\ge2$,
\begin{align*}
\mathbb{E}\bigl[1-\max(B_{p-1},x)^{\alpha}\bigr]
&=1-x^{\alpha}\,x^{p-1}-\tfrac{p-1}{p-1+\alpha}\bigl(1-x^{p-1+\alpha}\bigr)\\
&=\tfrac{\alpha}{p-1+\alpha}\bigl(1-x^{p-1+\alpha}\bigr),
\end{align*}
and for $p=1$ the factor is $1-x^{\alpha}$, which is the same
expression. This is the stated $\varphi_\alpha(p)$. Setting $x=0$ removes the
admissibility conditions and gives the rank distribution
\[
\theta_\alpha(p)=\frac{\alpha}{p-1+\alpha}\prod_{i=p}^{k-1}\frac{i}{i+\alpha}
=\frac{\alpha\,\Gamma(p-1+\alpha)\,\Gamma(k)}{\Gamma(p)\,\Gamma(k+\alpha)} .
\]
The mixture statement is linearity in the kind of the insertion.
\end{proof}

\paragraph{Consequences}
(1)~For $\alpha=1$, $\theta_1(p)=1/k$ and
$\varphi_1(p)=\frac1k\prod_{i=p}^{k}(1-x^i)$, so
$P_1(c)=1-\frac1k\sum_{p=1}^{k}\prod_{i=p}^{k}(1-x^i)
=\frac1k\sum_{i=1}^{k} i\,x^i+O(x^{3})=\frac{x}{k}+O(x^2)$, because
each product differs from $1-\sum_{i\ge p}x^i$ by terms of order
$x^{2p+1}$ or smaller.
(2)~Conditional on rank one, the insertion violates the first condition
of Theorem~5.1 with probability
$\Pr[\omega<x\mu_1\mid \omega<\mu_1]=x^{\alpha}=c^{-\Ds}$. (3)~As
$c\to\infty$, $1-\varphi_\alpha(p)/\theta_\alpha(p)=x^{p-1+\alpha}
+\sum_{i=p}^{k-1}x^{i+\alpha}+O(x^{1+2\alpha})$, whose leading term over
all ranks is $x^{\alpha}$ at $p=1$; therefore
$P_\alpha(c)= \theta_\alpha(1)\,c^{-\Ds}\,(1+o(1))$ and a mixture behaves as
$\sigma\,\theta_\alpha(1)\,c^{-\Ds}+(1-\sigma)\,c^{-D}/k$. (4)~The formulas remain
valid for $\alpha>1$, which describes writes that avoid the query, since the density
$\alpha\omega^{\alpha-1}/\mu_k^{\alpha}$ then grows with $\omega$.

\paragraph{Background rate}
We model a background write as exchangeable with the stored items, so
its rate needs no fitted parameter.
\begin{lemma}
\label{lem:exch}
Let $v$ be exchangeable with the $N$ items of $S$. Then $v$ is among
the $k$ nearest items of $S\cup\{v\}$ to a fixed query with probability
$k/(N+1)$, and given that it is, its rank is uniform on
$\{1,\dots,k\}$.
\end{lemma}
\begin{proof}
By exchangeability the rank of $\dist(q,v)$ among the $N+1$ distances
is uniform on $\{1,\dots,N+1\}$.
\end{proof}
Summing over the $n$ background writes that overlap a read gives the
expected number $n\,k/(N+1)$ of background in-ball insertions, with $N$
the store size at the read; no parameter is fitted.

\paragraph{Semi-parametric form}
For a read with neighbor distances $r_1,\dots,r_k$, a sibling insertion
has $(t/r_k)^{\Ds}$ uniform on $[0,1]$ given that it is in the ball, so
by Corollary~5.2 it conflicts with probability
$\Pr[t<\Rc(q)]=(\Rc(q)/r_k)^{\Ds}$. For a background insertion,
exchangeability says that, given the configuration of the nearest
items after the insertion, the inserted item is equally likely to be
any of them. We estimate the resulting conflict probability by
leave-one-out on the read's own neighbors, whose configuration has the
same law up to one item ($N$ stored items against $N+1$): designate the $p$-th of the read's $k+1$ nearest neighbors,
at distances $r_1<\dots<r_{k+1}$, as the inserted one, for $p\le k$ (an in-ball insertion). The held answer
would then be $r_1,\dots,r_{p-1},r_{p+1},\dots,r_{k+1}$, which is
inadmissible if and only if $r_{i+1}>c\,r_i$ for some $p\le i\le k$.
Averaging this indicator over $p=1,\dots,k$, which Lemma~\ref{lem:exch} makes equally likely, and over reads gives the background
prediction. With sibling shares $\sigma_j$
and sibling dimensions $\Ds^{(j)}$, one per sibling scale $j$, the predicted conflict probability of one
in-ball insertion at a read is
\[
\textstyle\sum_j\sigma_j\,(\Rc/r_k)^{\Ds^{(j)}}+\bigl(1-\sum_j\sigma_j\bigr)\cdot\text{(leave-one-out)}.
\]
For reads with several in-ball insertions we draw the positions of the
observed number of insertions from the same mixture around the read's
actual neighbor list and apply Lemma~\ref{lem:general}.

\subsection{Estimation}
\label{app:est}

All quantities are computed exact-to-exact from the logs of
instrumented runs: the auditor reconstructs, for the first read of
every task, the exact neighbor distances in the state at the read and
at the first validation, and the insertions committed in between.

\paragraph{Sibling dimensions}
For every in-ball insertion we record $y=\ln(r_k/t)$. Under the sibling
model $y$ is exponential with rate $\Ds$. The background share
$w_{\mathrm{bg}}$ is fixed at the ratio of the first-principles count
$\sum n\,k/(N+1)$ to the observed count, and the background density of
$y$ is the empirical law of $\ln(r_{k+1}/r_i)$, $i\le k$, over all
reads, smoothed by a three-component exponential mixture and then
frozen. We maximize the likelihood of the mixture
$w_{\mathrm{bg}}f_{\mathrm{bg}}(y)+\sum_j\sigma_j \Ds^{(j)} e^{-\Ds^{(j)} y}$ over one
or two sibling scales. Intervals come from a bootstrap that resamples
reads; goodness of fit is a Kolmogorov--Smirnov (KS) distance with a
parametric bootstrap. Conflict outcomes and ranks are never used for
fitting.

\paragraph{Lag kernel}
For a lag bucket $b$, $\hat h_b$ is the number of in-ball insertions
whose source item lies at a lag in $b$ behind or ahead of the reading
item in the same stream, divided by the number of writes at such lags
that overlapped the read, and $\hat D_b$ is the reciprocal of the mean
of $y$ over these insertions. The static estimator replaces the
concurrent run by the stream itself: for item $j$ it takes the exact
top-$k$ ball in the state that contains the initial memory and the
items before $j$, and counts the items $j+\ell$, $\ell\in b$, that fall
into it. A prediction for a run uses the run's overlap counts $o_b$ and
the kernel: $\nu=\sum_b o_b\hat h_b$ in-ball insertions per read, and
the mixture of sibling laws with weights $o_b\hat h_b/\nu$ and
dimensions $\hat D_b$ in the semi-parametric form, one scale per lag bucket.

%% ====================================================================
\section{PostgreSQL: Details of \S3}
\label{app:pg}

\subsection{Background: serializable snapshot isolation in PostgreSQL}
\label{app:ssi-bg}

PostgreSQL implements its serializable level as serializable snapshot
isolation (SSI) \cite{DBLP:conf/sigmod/CahillRF08,DBLP:journals/pvldb/PortsG12}.
Every transaction reads from a snapshot, and SSI tracks
\emph{rw-antidependencies}: $T_r\to T_w$ when $T_r$ reads a version
that $T_w$ later replaces, or when $T_w$ inserts a row that a predicate
read of $T_r$ would have returned (a phantom). Reads are recorded by
predicate locks (SIREAD locks), which never block: a read takes them on
the tuples it returns and on the index pages, or the whole index
relation, that its scan visits, and every write checks them through
\texttt{CheckForSerializableConflictIn}. Every cycle that snapshot
isolation admits contains two consecutive rw-antidependencies between
concurrent transactions, a \emph{dangerous structure}
\cite{DBLP:journals/tods/FeketeLOOS05}; SSI aborts a transaction of
every dangerous structure it detects, so it may also abort executions
that are serializable. An index access method declares through
\texttt{ampredlocks} whether it takes page-level predicate locks itself.
B-tree, hash, GiST, and GIN do; for a method that does not, a scan
locks the whole index relation, and \texttt{index\_insert} checks the
whole relation before it calls the method's insert routine. Both
properties matter in \S3.2: the relation-level lock makes every
similarity read conflict with every insertion (coarseness), and the
order of the check and of the insertion opens the window described
next (unsoundness).

\subsection{Code path and the fix}

If the handler of an index access method does not set
\texttt{ampredlocks}, \texttt{index\_insert} calls
\texttt{CheckForSerializableConflictIn} on the whole relation and then
the method's insert routine. For pgvector's HNSW this routine,
\texttt{hnswinsert}, searches the graph for the neighbors of the new
element and then links it. A scan by another serializable transaction
takes its relation-level predicate lock when it starts. If it starts
after the writer's check, the writer has not seen its lock; if it
finishes before the element is linked, it has not seen the element.
Figure~\ref{lst:patch} repeats the check after the element is
reachable. It is the smallest change that closes the window; it does
not make the lock finer than the whole index relation. We reported the window and this change to the pgvector
maintainers \cite{pgvectorissue}.

\begin{figure}[t]
\begin{Verbatim}[fontsize=\scriptsize,frame=single]
--- a/src/hnswinsert.c
+++ b/src/hnswinsert.c
@@ -3,6 +3,7 @@
 #include "hnsw.h"
+#include "storage/predicate.h"
@@ -793,6 +794,13 @@ hnswinsert(Relation index, ...
 	/* Insert tuple */
 	HnswInsertTuple(index, values, isnull, heap_tid);

+	/* repeat the SSI conflict-in check after the
+	 * element is linked into the graph */
+	CheckForSerializableConflictIn(index,
+	    (ItemPointer) NULL, InvalidBlockNumber);
+
 	/* Delete memory context */
\end{Verbatim}
\caption{The change to pgvector (abridged diff) that closes the window
  of \S3.2 in \texttt{hnswinsert}.}
\label{lst:patch}
\end{figure}

\subsection{Tests}

\autotab{tab_ssi_soundness}{tab:app-ssi}
  {Soundness tests for serializable isolation over index access
  methods (generated from the test logs). In the ivfflat control, the
  reader's probes miss the committed row in some trials: an index miss,
  not a race.}

\autofig{fig_s_06_ssi_dose}{fig:app-dose}
  {Share of transaction pairs that store both rows against the mean insertion time, with intervals, higher for the slower configurations than for the default; the advisory-lock and sequential-scan controls at zero, and the patched library at \res[N]{A6.patched.dup} of \res[0{,}000]{A6.patched.pairs} pairs.}
  {Share of tightly timed serializable pairs that store both
  rows, against the mean duration of an HNSW insertion (varied through
  $M$ and $\mathit{ef\_construction}$), with 95\% Wilson intervals; hollow marks: the controls.}

\autofig{fig_s_07_ssi_cost}{fig:app-ssicost}
  {Bars of aborts per task and model calls per task for each protocol at several concurrency levels, with a control on a scalar key.}
  {Aborts and model calls per task under serializable isolation
  and its alternatives, by concurrency level (light to dark: 4, 16, 32
  agents; hatched: 32; ``global lock, whole txn.'' is the lock across
  the thinking step); the printed numbers are the values at 16 agents, as in
  Table~\ref{M-tab:ssi} of the paper.}

Table~\ref{tab:app-ssi} lists all tests: the predicate locks that one
read takes per access method, the deterministic reproduction with a
debugger pause, the stress test without a debugger, and its controls.
Figure~\ref{fig:app-dose} shows the share against the mean duration
of an insertion: every configuration with slower insertions than the
default lets a tightly timed reader fall into the window more often,
most often (\ra{A6.grid.max.rate}\%) at a mean insertion time of \ra{A6.grid.max.ms}\,ms. The same pause on other
access methods without predicate locking (SP-GiST, BRIN) produces a
double commit; on B-tree, hash, and GIN, which take page-level
predicate locks inside the insertion, one of the two transactions
aborts; GiST, which does lock the pages a search visits, also commits
both rows (\S\ref{app:gist}), and the pause is never reached in bloom. Figure~\ref{fig:app-ssicost} complements Table~\ref{M-tab:ssi} of the paper with the
other concurrency levels and the control on a scalar key indexed by a B-tree.

\subsection{GiST}
\label{app:gist}

Our reading of the observed behavior is the following. GiST takes
predicate locks on the pages that a search visits and checks an
insertion against the locks on the pages it modifies. A search for a
point outside every bounding box of an internal page is pruned there
and locks only that page. The first insertion of such a point enlarges
a bounding box on the internal page, modifies it, and is checked
against the lock. A second insertion of the same point now descends
into the enlarged box and modifies only a leaf page, which no search
locked. If two transactions search before either inserts, both find
nothing, the first insertion conflicts with the second transaction's
lock (one dependency), and the second insertion conflicts with nothing,
so no dangerous structure forms. The protocol of Kornacker et
al.~\cite{DBLP:conf/sigmod/KornackerMH97} attaches the search predicate to the node and
propagates it to a child when the child's bounding predicate grows; the
implementation moves predicate locks on page splits only. We observed
double commits in \ri{A7.out.eq} of \ri{A7.trials} trials for equality
and \ri{A7.out.range} for range searches outside all data,
\ri{A7.in.eq} of \ri{A7.trials} for points inside the populated region,
\ri{A7.out.knn} for nearest-neighbor ordering, and \ri{A7.out.btree}
for the same test on a B-tree. We have not
traced the mechanism in the source line by line, and we tested one
major version.

%% ====================================================================
\FloatBarrier
\section{Conflict Law: Details of \S\S5 and~6.5}
\label{app:law}

\autofig{fig_t_11_characterization}{fig:app-char}
  {A number line of the neighbor distances of one real read on a logarithmic axis, with the interval of insertion distances that conflict shaded for each rank.}
  {Theorem~5.1 on the neighbor list of one real read. (a)~For each
  rank, the distances at which an insertion breaks admissibility;
  (b)~the conflict region over all tolerances.}

\autofig{fig_t_13_lid_sensitivity}{fig:app-lid}
  {Error of the closed form per configuration for three estimators of the intrinsic dimension.}
  {Sensitivity of the closed form to the estimator of the
  intrinsic dimension. Black ticks: the semi-parametric form (two
  sibling scales on AG News, one elsewhere).}

\autofig{fig_t_08_controls}{fig:app-controls}
  {Bars of in-ball insertions per read for the natural order, block-shuffled order, shuffled order, and an exchangeable split, each with a marker for the first-principles background rate.}
  {Controls on LoCoMo (one team of 80 agents, unlike L8) and A64 as the order of the stream
  is destroyed step by step (blocks of 64 items shuffled, all items shuffled, exchangeable
  split). (a)~In-ball insertions per read (log scale); black ticks:
  first-principles background rate. (b)~Conflict probability after one
  in-ball insertion at $c=1.1$; black ticks: exchangeable-only
  prediction. (c)~Share of all
  reads with a stale answer, for each order of the stream.}

Besides the configurations of \S6.1, the law runs include A256 (AG News, 256 agents, 12{,}000 writes) and A64-long (AG News, 64 agents, 30{,}000 writes), shown in Figure~\ref{M-fig:lag} of the paper and below.

\paragraph{Closed form}
Figure~\ref{fig:app-heat} maps Theorem~5.3. For a pure sibling stream,
the $c-1$ needed for a given conflict probability grows as the sibling
dimension falls; background writes
flatten this dependence where conflicts are likely.
Figure~\ref{fig:app-char} applies Theorem~5.1 to the neighbor list of
one real read.
Figure~\ref{fig:app-mc} compares the closed form with simulation of the
model on \ri{D1.cells} parameter settings; the largest standardized
deviation is \ra{D1.maxz}, and \ra{D1.above2}\% of the tested
combinations of setting and seed exceed two standard errors of the
Monte Carlo estimate of the same parameter setting and seed.

\paragraph{Where insertions land}
Figure~\ref{fig:app-landscape} is the data-driven view of
Corollary~5.2: for every read of a run we compute $\Rc(q)/r_k$ and
average the indicator that an insertion at normalized distance
$t/r_k$ conflicts. The marginal on top shows where real in-ball
insertions land, split into the first-principles background and the
excess over it. Background insertions land near the rim, where almost
no tolerance is needed; the excess lies deeper.
Figure~\ref{fig:app-lawall} repeats Figure~\ref{M-fig:law} of the
paper for every configuration, Figure~\ref{fig:app-ydens} shows the
fitted position densities, and Figure~\ref{fig:app-rank} the rank
distributions, unseen by the fit.

\paragraph{Ranks}
The position fit never sees the ranks of the in-ball insertions. Its
predicted share at rank one is within \ra{D3.rank1.maxdiff} points of the
observed share on the six configurations of Table~\ref{M-tab:law} of
the paper, far closer than uniform ranks (Figure~\ref{fig:app-rank}).
A chi-square test of the whole rank distribution, pooled over seeds and
corrected for the dependence of insertions within a burst, rejects the
prediction on these configurations except L4, so the law describes the
share at rank one, not the whole rank distribution.

\paragraph{All reads}
For reads with any number of in-ball insertions, the simulation of
\S\ref{app:law-proof} predicts the share whose held answer is
inadmissible at its first validation; the largest error over configurations and
tolerances is \ra{D3.e2e.maxerr} points (two sibling scales on AG News), against
\ra{D3.e2e.bg.maxerr} for background writes only. Stratified by the
number of in-ball insertions of the read, the errors are
\ra{D3.strata.m1.err}, \ra{D3.strata.m2.err}, and
\ra{D3.strata.m3.err} points for one, two, and three or more (largest
over the configurations with at least 30 such reads; this check uses
one sibling scale per configuration). The larger errors are on AG News,
which needs two scales (Table~\ref{M-tab:law}); on LoCoMo, which one
scale describes, the stratum of reads with one in-ball insertion errs least.

\paragraph{Tails and dimension estimators}
Figure~\ref{fig:app-tail} compares the log-log slope of the observed
conflict probability with the fitted sibling dimensions, and
Figure~\ref{fig:app-lid} shows how the error of the closed form depends
on the estimator of the intrinsic dimension (the semi-parametric form
needs none). Stratifying the reads by their local intrinsic dimension
raises the error of the closed form with one pooled $D$ to
\ra{D3.strata.lid.closed} points on average and that of the
semi-parametric form to \ra{D3.strata.lid.semi}; both rise by similar
amounts. On a synthetic store that mixes two manifolds of different
dimension, the closed form errs more than the semi-parametric form
(Figure~\ref{fig:app-hetero}).

\paragraph{Planted law and controls}
On synthetic streams the background items are exchangeable with the
stored ones and siblings are planted at chord distance
$R\,\xi^{1/\Ds}$ from a parent a few positions earlier, with $\xi$
uniform on $(0,1)$ and a radius $R$ above the $(k{+}1)$-th neighbor
distance of almost every read.
Figure~\ref{fig:app-planted} compares planted and recovered sibling
dimensions over \ri{D2.total} configurations and seeds (the figure
draws means over seeds); the planted value lies inside the 95\%
interval of its run in \ri{D2.inside} of them. On these synthetic
streams (not drawn), the pooled in-ball rate of background insertions is \rb{D2.bg.ratio} times the
first-principles rate, and at $k=50$ their rank-one share is
\ra{D2.bg.rank1}\% (95\% interval \ra{D2.bg.rank1.lo} to
\ra{D2.bg.rank1.hi}; uniform ranks give \ra{D2.bg.uniform}\%). Figure~\ref{fig:app-controls}
shows the controls of \S6.5 on real streams. In the exchangeable split
the rank-one share of in-ball insertions is \ra{D4.exch.rank1}\% (95\% interval
\ra{D4.exch.rank1.lo} to \ra{D4.exch.rank1.hi}, against \ra{D4.exch.uniform}\% for uniform ranks), and a
chi-square test of the whole rank distribution does not reject uniform
ranks ($p=\rb{D4.exch.chip}$).

\paragraph{Transfers}
Figure~\ref{fig:app-transfer} shows every prediction that was made
without fitting on its target: across concurrency levels with pooled
parameters and with the lag kernel, across $k$, across time, from the
static stream, and across collections. Figure~\ref{fig:app-kemb} shows
the fitted parameters for other values of $k$ and other embedding
models.

\autofig{fig_t_04_law_heatmap}{fig:app-heat}
  {Two filled contour maps of the closed-form conflict probability over tolerance and sibling dimension, for a pure sibling stream and for an even mixture with background writes, with the fitted sibling dimensions of the workloads in a strip on the right.}
  {Conflict probability of one in-ball insertion by Theorem~5.3
  over the tolerance and the sibling dimension, for two sibling shares.
  The strip on the right marks the fitted sibling dimension $\Ds$ of each workload.}

\autofig{fig_t_14_monte_carlo}{fig:app-mc}
  {Scatter of simulated against closed-form conflict probabilities on the diagonal, with a histogram of standardized deviations.}
  {Closed form of Theorem~5.3 against Monte Carlo simulation of
  the model, per setting and seed; hollow marks: too few simulated
  conflicts for a standardized deviation. (c)~Standardized deviations of
  the conflict probability, the rank-one share, and the conflict
  probability given rank one; the maximum and the share above two in the
  text refer to the conflict probability.}

\autofig{fig_t_02_conflict_landscape}{fig:app-landscape}
  {Heat maps of the probability that an insertion at a normalized distance from the query breaks admissibility, over that distance and the tolerance, for two streams, with the distribution of real insertions on top.}
  {Probability that an insertion at normalized distance
  $t/r_k$ breaks $c$-admissibility, averaged over the real neighbor
  lists of the runs of a data set, and where concurrent insertions land.
  Dotted lines: median position of the observed insertions and of the
  background law; the hatched last bin also holds the insertions with
  $1-u<0.004$. Right: conflict probability by tolerance for the observed
  insertions (solid) and the background law (dashed).}

\autofig{fig_t_03_law_vs_observed}{fig:app-lawall}
  {Small multiples, one per configuration, of the conflict probability against tolerance: observed points with intervals, the exchangeable prediction, the closed form, and the semi-parametric form with one and two scales.}
  {Conflict probability of reads with one in-ball insertion,
  all configurations (the paper shows two of them).}

\autofig{fig_t_10_y_density}{fig:app-ydens}
  {Histograms of the position of in-ball insertions per configuration with the fitted one-scale and two-scale sibling mixtures and the background shape.}
  {Positions $y=\ln(r_k/t)$ of in-ball insertions with the fitted
  mixtures and the background shape. Printed: KS distance (mean over
  seeds; the paper quotes the largest single run) with its parametric bootstrap $p$ (range over seeds;
  $p\leq 0.005$ is the smallest value that 200 bootstrap draws can give).}

\autofig{fig_t_09_rank_distribution}{fig:app-rank}
  {Bar charts of the old rank of in-ball insertions per configuration: observed, predicted by the fitted law, and uniform.}
  {Old rank of in-ball insertions: observed, predicted from the
  position fit with one and with two sibling scales (never fitted on
  ranks; the text uses one scale on LoCoMo and two on AG News), and the
  uniform distribution of exchangeable writes.}

\autofig{fig_t_12_tail_exponents}{fig:app-tail}
  {Observed log-log slopes of the conflict probability per configuration with intervals, next to the fitted sibling dimensions and the slope of the exchangeable model.}
  {Tail exponents of the conflict probability against the fitted
  one-scale sibling dimensions.}

\autofig{fig_t_15_heterogeneous_dimension}{fig:app-hetero}
  {Prediction errors of the closed form and the semi-parametric form on synthetic data that mixes two manifolds of different dimension, overall and by stratum of local dimension.}
  {Heterogeneous intrinsic dimension: (a)~closed form with one pooled
  $D$ against the semi-parametric form on synthetic stores; (b)~errors
  inside strata of local dimension, synthetic and real data.}

\autofig{fig_t_07_planted_recovery}{fig:app-planted}
  {Recovered sibling dimension (mean over runs) against planted sibling dimension on the diagonal, and observed against first-principles background rate on the diagonal.}
  {(a)~Recovered against planted sibling dimension on synthetic
  streams; (b)~observed against first-principles background rate on
  the exchangeable controls of real data.}

\autofig{fig_t_06_transfers}{fig:app-transfer}
  {Panels of observed and predicted conflict curves of target runs for several transfer modes, with the maximum error of each.}
  {Predictions without fitting on the target, by transfer mode (marker
  shape: kind of transfer). The target's own fit uses the scales of
  its row: two on AG News for the static stream, those of the transfer
  otherwise, so one target can show two own-fit errors, one per row in which it appears.}

\autofig{fig_t_16_k_and_embedders}{fig:app-kemb}
  {Fitted sibling dimension, sibling share and conflict curves against the retrieval depth, per data set and embedding model.}
  {Other retrieval depths and embedding models. (a)~In-ball
  insertions per read with the first-principles background; (b,~c)~the
  fitted sibling share and dimension; (d)~observed (marks) and predicted
  (lines) conflict probability at $c=1.1$, observed marks of the models
  offset slightly along $k$. Color: data set; line style and marker:
  embedding model (WordLlama: the data set's filled marker; the others:
  hollow markers of their own).}

%% ====================================================================
\FloatBarrier
\section{Additional Experiments}
\label{app:exp}

\subsection{The race (\S3.1)}
\label{app:exp-race}

\autofig{fig_s_02_race_spread}{fig:app-jitter}
  {Duplicates per 100 facts against the arrival spread of the copies in units of the thinking time, falling from high values to near zero.}
  {Unprotected check-then-insert on pgvector: duplicates against
  how tightly co-observations arrive.}

\autofig{fig_s_04_race_realtext}{fig:app-real}
  {Duplicates for real duplicate question clusters and duplicate writes on replayed streams, by protocol.}
  {Real paraphrases and replayed streams.}

Figure~\ref{fig:app-race} shows the race per store at every concurrency
level, Figure~\ref{fig:app-jitter} its dependence on the arrival spread
of the copies, and Figure~\ref{fig:app-factors} the remaining factors
one at a time: thinking time, fan-out, threshold, retrieval depth,
topic locality, store size, and concurrency. Of the timing and load
parameters the arrival spread relative to the thinking time matters
most; concurrency and the others shift the rate far less; the fan-out sets how many copies there are to collide, and the
threshold does not matter while it exceeds the distance between copies.
Figure~\ref{fig:app-real} repeats the measurement with real duplicate
questions and with replayed streams, where the auditor counts
insertions made although an item within the threshold existed.
Figure~\ref{fig:app-anomalies} breaks the anomalies down by workload
and protocol. Through Mem0 (its default pipeline with Chroma as
the store, one shared user, a target of eight calls in flight),
about \ri{A9.adds} concurrent add calls per run for \ri{A9.events}
facts (three runs) stored \ra{A9.concurrent.dup} duplicates per 100 facts (serial:
\ra{A9.serial.dup}; with a lock around the call: \ra{A9.lock.dup}).
Writing first and verifying afterward is no remedy either: on the
workload of Table~\ref{M-tab:race} of the paper, deleting the own copy when a near duplicate with a
smaller identifier is found leaves \ra{A11.wvid.dup} duplicates per 100
facts, and deleting it whenever one is found leaves
\ra{A11.wvany.dup} and loses \ra{A11.wvany.lost} facts per 100, whose
copies all deleted themselves.

\autofig{fig_s_01_race_stores}{fig:app-race}
  {Grouped bars of duplicates per 100 facts for each store and protocol, one panel per concurrency level.}
  {Duplicates per 100 facts by store and protocol.}

\autofig{fig_s_03_race_factors}{fig:app-factors}
  {Small multiples of duplicates per 100 facts against each workload factor.}
  {The race against each workload factor, one at a time
  (unprotected check-then-insert on pgvector; steps equally spaced; the
  fan-out panel has its own y scale).}

\autofigtop{fig_s_34_anomalies}{fig:app-anomalies}
  {Annotated table of anomaly rates by workload and protocol.}
  {Anomalies by workload (adds, upserts, mixed operations) and
  protocol; counts per 100 tasks unless marked \%, shading relative to
  the column maximum within a panel (0.00: a mean below 0.005 that is
  not zero; gray: a cost, not an anomaly). The $c$-admissible rows validate
  with an approximate read, so their answers are bounded by $c\rho$
  rather than $c$; in the other rows, answers above the ratio come from
  the index itself.}

\subsection{Correctness checks (\S4)}
\label{app:exp-correct}

\autotab{tab_safety_all}{tab:app-safety}
  {Audited commits and violations of the safety bound, per
  experiment, for the protocol as specified and for variants weakened
  on purpose (generated).}

\autofig{fig_s_25_separation}{fig:app-sep}
  {Smallest stored pairwise distance divided by the threshold against the tolerance, with the curve one over c.}
  {Smallest stored pairwise distance against the tolerance, with
  the bound $\tau/c$ of Proposition~4.2 for certified runs; with approximate validation reads the bound is
  $\tau/(c\rho)$~(\S\ref{app:sep}). Caps: 95\% Student-$t$ intervals
  of the mean over three seeds, cut at 0; series are offset slightly
  along $c$ so that coinciding marks of different series stay visible.}

Table~\ref{tab:app-safety} lists, per experiment, the number of commits
on which the inequality of Theorem~4.3 was checked and the number of
violations. Runs of the protocol as specified are separated from runs
of variants that we weakened on purpose (no existence check although
items are deleted, no commit lock or a lock narrower than the read
scope, a snapshot taken before the lock is granted). Violations in the
second group come from the runs without the existence check; they show
that the check is not vacuous. The variants of the lock and of the
snapshot pass it, since the inequality holds between the held answer
and the validation read they take; when the workload exposes them,
they fail through duplicate writes
and fresh answers above the ratio (\S\ref{app:audit}). Figure~\ref{fig:app-sep} shows the smallest stored pairwise
distance against the bound $\tau/c$ of Proposition~4.2, including a
synthetic stream whose siblings are planted around the threshold so
that the bound is approached.
The auditor computes the ratio of a held answer against the state at
the start of its task. With several reads per task or with deletions, a
later read can return an item inserted after that state or miss one
deleted after it, so on A64multi and A64mix the share of held answers
with a ratio above $c$ is not zero even in certified runs. The safety
inequality, which compares the held answer with the fresh answer in the
commit state, has no violation in these runs. On LongMemEval and the
question stream, which contain exact copies, a few certified answers
hold an item at distance zero; the store returns that distance as zero
and the auditor recomputes it from the vectors as the float32
resolution of the chord distance, so these exact answers show a large
ratio. Every certified answer above $c$ outside A64multi and A64mix is
of this kind.

\subsection{Cost and savings (\S\S6.2 to~6.4)}
\label{app:exp-cost}

\autotab{tab_cost_all}{tab:app-cost}
  {Cost of every validation predicate on every configuration, all on
  pgvector (generated; eight seeds on L2, L4, and A64, three on the
  others). On MS MARCO, whose base of one million passages makes each
  run slow, only the predicates none and exact and the tolerances
  $c\in\{1.05, 1.1, 1.2\}$ were run.}

\autofig{fig_s_14_rho_ccdf}{fig:app-rho}
  {Complementary distribution of the approximation ratio of held answers per predicate, with the distribution of fresh answers and vertical lines at each tolerance.}
  {Approximation ratio of held answers per predicate
  (complementary cumulative distribution function, CCDF), and of fresh answers for comparison.}

\autofig{fig_s_24_choosing_c}{fig:app-choose}
  {Extra calls and differing decisions per decision function against the tolerance on a fine grid, with reference lines at quantiles of the index's approximation ratio.}
  {Choosing $c$: extra calls and differing decisions on a fine grid of tolerances.
  Vertical lines: 90th and 99th percentiles (where above 1; on LoCoMo the 90th is 1) over all tasks of the
  index's own approximation ratio (worst validation read per task),
  averaged over the runs.}

\autotab{tab_llm}{tab:app-llm}
  {Two language models, \res[two language models]{env.llm}, decide on L4, L8, and
  LongMemEval, five seeds per cell; changed decisions from a replay of
  \res[0{,}000]{C3.replay.n} tasks per run, F1 and judge accuracy of the downstream questions on the final memory of
  every seed (generated; the questions exist for L4 and LongMemEval, so L8 has none).}

Table~\ref{tab:app-cost} extends Table~\ref{M-tab:cost} of the paper to all predicates and
configurations, and Figure~\ref{fig:app-pareto} extends the frontier to
all workloads. Figure~\ref{fig:app-rho} shows the distribution of the
approximation ratio of held answers per predicate against that of
fresh answers; Figure~\ref{fig:app-choose} sweeps $c$ on a fine grid (over the grid, extra calls fall and differing decisions do not fall as $c$ grows)
with the quantiles of the index's own ratio as reference;
Figure~\ref{fig:app-budget} varies the index and its search budget.
Table~\ref{tab:app-llm} is the generated version of Table~\ref{M-tab:llm} of the paper for both
models and all three workloads (L4, L8, and LongMemEval).

\autofig{fig_s_10_pareto_all}{fig:app-pareto}
  {The trade-off between model calls saved and differing decisions for every workload, one row per workload and one panel per decision function.}
  {Calls saved against differing decisions, all workloads (rows: data
  set and number of agents). The index floor counts all transactions,
  the points only accepted decisions, so a point can lie below
  it. MS MARCO passages have no tag labels, so the majority-tag panel shows no points for MS MARCO.}

\autofig{fig_s_15_contract_vs_budget}{fig:app-budget}
  {Share of answers whose ratio exceeds the tolerance against the search budget, for held and fresh answers, per index and validation variant.}
  {Answers above the tolerance ($c=1.1$) against the search budget
  of the index.}

\FloatBarrier
\subsection{Performance (\S6.6)}
\label{app:exp-perf}

\autotab{tab_opt_vs_pess}{tab:app-optpess}
  {Optimistic against pessimistic: all six workloads at every
  tolerance, three seeds; ``validation'' is certified validation
  (exact validation read), and both designs run in the client library
  on the in-memory store (generated).}

\autofig{fig_s_11_scaling_all}{fig:app-scaleall}
  {Log-log panels of throughput and 99th percentile latency against the number of agents, one row per thinking time, one line per protocol.}
  {Throughput and tail latency against the number of agents at every
  thinking time.}

\autofig{fig_s_16_critical_section}{fig:app-cs}
  {Stacked bars of the time in the critical section (validation read, predicate, write), and the lock wait in a separate panel.}
  {Time in the critical section by its parts, for retrieval depths
  $k$ and search budgets of the validation read, and with an exact
  validation read (AG News, pgvector, $c=1.1$). (a)~Parts of the
  critical section, upper bar 16 and lower bar 128 agents; (b)~lock wait,
  light 16 and dark 128 agents.}

\autofig{fig_s_13_think_time}{fig:app-think}
  {Extra calls against the mean thinking time for five distributions.}
  {Thinking-time distributions.}

\autofig{fig_s_12_open_loop}{fig:app-open}
  {Latency and completed rate against the arrival rate for Poisson and bursty arrivals.}
  {Open-loop arrivals. Bursty marks are drawn slightly to the right of
  their arrival rate so that they do not hide the Poisson marks. In (b)
  every protocol completes nearly the offered rate, so the marks coincide.}

Figure~\ref{fig:app-optpess} and Table~\ref{tab:app-optpess} compare
validation with ball predicate locks on every workload at every
tolerance; on the LoCoMo and AG News streams, tolerance reduces the extra calls
of both designs, and on the fan-out workload, where the tasks of a team
read the same neighborhood, the extra calls of both stay level.
Figure~\ref{fig:app-rounds} gives the distribution of rounds per task
in the worst cases. Figure~\ref{fig:app-scaleall} shows throughput and tail latency
against the number of agents at both thinking times: from 16 agents
on, $c=1.1$ and re-execution keep about \ri{C5.a64.c110.rel}\% and \ri{C5.a64.exact.rel}\% of the
throughput without isolation (\ri{C5.a256.c110.rel}\% and \ri{C5.a256.exact.rel}\% at 256 agents, where
$c=1.1$ holds the commit lock \ri{C5.a256.c110.held}\% of the time, \ri{C5.t2.a256.c110.held}\% with thinking times
of 2\,s); certified validation, whose exact scan holds the lock
\ra{C1.A64.cert110.cs}\,ms per task, saturates it from 16 agents on at 0.2\,s
(\ri{C5.a16.cert110.held}\%) and stays near \ri{C5.a256.cert110.tps} tasks per second, and at 2\,s it
saturates only from 128 agents on (\ri{C5.t2.a64.cert110.held}\% at 64, \ri{C5.t2.a256.cert110.held}\% at 256);
ball locks fall further behind as their deadlocks grow with the agents;
Figure~\ref{fig:app-cs} breaks the critical section down,
Figure~\ref{fig:app-think} varies the distribution of the thinking
time, Figure~\ref{fig:app-open} replaces the closed loop by open-loop
arrivals, and Figure~\ref{fig:app-scale} varies the size of the memory
and of the run. The long run keeps only its
summary rows, so its time series is not shown.
Throughput times mean latency equals the number of clients within a
factor of 0.8 to 1.1 in every closed-loop run, so no client waits
outside the measured latency; with one process per agent the smallest
ratio is \rb{F4.little.min}.

\autofig{fig_s_18_opt_vs_pess}{fig:app-optpess}
  {Scatter of extra model calls against lock wait for validation and ball locks over tolerances.}
  {Validation against ball predicate locks, all six workloads. Each
  path is labeled at its ends ($c=1$ and $c=1.5$); numbers in
  parentheses are deadlocks per 100 tasks; horizontal intervals are
  clipped at 0.}

\autofig{fig_s_19_rounds_ccdf}{fig:app-rounds}
  {Complementary distribution of the number of rounds per task under worst-case workloads.}
  {Rounds per task under worst-case workloads (CCDF), all with
  heavy-tailed thinking: AG News with 256 agents, LoCoMo with 8, and hot
  topics on the fan-out workload. Solid: a cap of 30 rounds; dashed,
  hollow marks: at most 3 rounds, then commit under the lock. On the
  hot-topic workloads, $c=1.2$ finishes every task within \ri{C11.fan16.c120.maxrounds} rounds and
  re-execution needs up to \ri{C11.fan16.exact.maxrounds}.}

\autofig{fig_s_23_scale}{fig:app-scale}
  {Extra calls, validation read time and throughput against the size of the memory.}
  {Size of the memory and of the run: (a)~extra calls per 100 tasks, (b)~mean time of the validation
  read (validating protocols only), (c)~throughput (log scale). The size
  sweep uses 64 agents and 6{,}000 tasks; the two right-most cases start
  from 60k items, one with 512 agents (20{,}000 tasks), one with 50{,}000 tasks.}

\FloatBarrier
\subsection{Ablations and robustness (\S6.7)}
\label{app:exp-abl}

\autotab{tab_ablation}{tab:app-abl}
  {Ablations (generated).}

\autofig{fig_s_30_retry}{fig:app-retry}
  {Extra calls, capped transactions and tail latency per strategy after a failed validation and round limit.}
  {Strategy after a failed validation, by round limit (limits evenly
  spaced; the axis is broken before 30).}

\autofig{fig_s_17_validation_budget}{fig:app-vbudget}
  {Validation read time, answers above the tolerance and throughput against the search budget of the validation read.}
  {Search budget of the validation read: (a)~read time (log scale),
  (b)~answers above~$c$, (c)~throughput.}

\autofig{fig_s_31_isolation_phases}{fig:app-phases}
  {Duplicate writes and extra calls per isolation level of the read phase and of the critical section.}
  {Isolation levels of the two phases on PostgreSQL.}

\autofig{fig_s_32_per_decision}{fig:app-perdec}
  {Extra calls, differing actions and differing link targets for uniform and per-operation tolerances.}
  {Per-operation tolerance.}

\autofig{fig_s_21_k_tau}{fig:app-ktau}
  {Extra calls against the retrieval depth and against the threshold.}
  {Extra calls against the retrieval depth and the threshold; each sweep varies one parameter with
  the other at its default (dotted line: $k=10$, $\tau=0.1$).}

\autofig{fig_s_22_read_mix}{fig:app-readmix}
  {Extra calls, throughput and the 99th percentile of the approximation ratio against the share of read-only tasks.}
  {Share of read-only tasks (AG News, 64 agents).}

\autofig{fig_s_27_stores}{fig:app-stores}
  {Extra calls, duplicate writes and critical-section time per store.}
  {Stores: (a)~extra calls, (b)~duplicate writes, (c)~mean time in the
  critical section (log scale; without isolation, the time of the
  write itself).}

\autofig{fig_s_28_failures}{fig:app-fail}
  {Throughput, tail latency and duplicate writes against the probability of a client failure in the critical section.}
  {Client failures inside the critical section.}

\autotab{tab_portability}{tab:app-port}
  {Portability across stores (generated).}

Table~\ref{tab:app-abl} is the generated ablation table with all
blocks, including the strategy after a failed validation and the
per-operation tolerance. The ablations are shown per figure: the
components of the predicate (Figure~\ref{fig:app-pred}), the scope of
the commit lock (Figure~\ref{fig:app-lock}), the strategy after a
failed validation (Figure~\ref{fig:app-retry}), the search budget of
the validation read (Figure~\ref{fig:app-vbudget}), the isolation
levels of the two phases (Figure~\ref{fig:app-phases}), and the
per-operation tolerance (Figure~\ref{fig:app-perdec}). The robustness
experiments cover several reads per transaction
(Figure~\ref{fig:app-multi}), the retrieval depth and the threshold
(Figure~\ref{fig:app-ktau}), the share of read-only tasks
(Figure~\ref{fig:app-readmix}), the stores (Figure~\ref{fig:app-stores}),
and client failures (Figure~\ref{fig:app-fail}). Table~\ref{tab:app-port} lists the stores.

\autofig{fig_s_33_predicates}{fig:app-pred}
  {Extra calls, duplicate writes, answers above the tolerance and differing decisions per predicate.}
  {Predicate components.}

\autofig{fig_s_29_lock_scope}{fig:app-lock}
  {Lock wait, throughput and duplicate writes per lock scope.}
  {Scope of the commit lock (top: re-execution validation; bottom:
  $c=1.1$; fill and hatch: lock scope). LoCoMo has one team per
  conversation: 10 teams of 2, 8 or 32 agents.}

\autofig{fig_s_20_multi_read}{fig:app-multi}
  {Extra calls and critical-section time against the number of reads per transaction.}
  {Transactions with several reads.}

%% ====================================================================
\section{Validating the Auditor}
\label{app:audit}

\autofig{fig_s_35_broken_variants}{fig:app-broken}
  {Duplicate writes and share of fresh answers above ratio 1.1 for the correct protocol and three broken variants.}
  {Deliberately broken variants of the protocol, as seen by the
  auditor: (a)~duplicate writes per 100 tasks (symmetric log scale;
  means on the bars), (b)~fresh answers above ratio 1.1 (\%).
  Duplicates counted from the validation reads equal those in (a) in
  every run.}

The auditor is the instrument behind every quality number, so we
tested whether it detects protocols that are wrong.

\paragraph{Broken variants}
We ran three variants that look correct: validation without mutual
exclusion; one commit lock per team on a store shared by all teams; and,
on PostgreSQL, a critical section at repeatable read, whose snapshot is
taken before the advisory lock is granted. Figure~\ref{fig:app-broken}
shows the outcome. Without a lock, co-arriving copies validate
concurrently and are stored twice (\rb{B5a.fan.nolock.dupw} duplicates per 100 tasks on the
fan-out workload); the correct protocol writes only the
duplicates that the index itself misses (\rb{B5a.fan.correct.dupw} per 100 tasks on the
fan-out workload, its index floor). One lock per team is
indistinguishable from the correct protocol on the fan-out workload,
which has one team, and on LoCoMo, whose teams write about different
conversations. On the news stream split into teams, which write about
the same things, it stores \rb{B5a.A64rr8.group.dupw} duplicates per 100 tasks
against \rb{B5a.A64rr8.correct.dupw}, a difference that three seeds do not
resolve (the ten seeds of the lock-scope ablation do, Table~\ref{tab:app-abl}): the variant is unsafe by construction, but this workload
rarely exposes it. The repeatable-read variant is instructive: the
inequality of Theorem~4.3 between the held answer and the
fresh answer still holds, because the fresh answer is itself
stale; the auditor exposes it through the ratio of the fresh
answer measured in the commit state, which exceeds 1.1 for
\ra{B5a.rrcs.excess.min} to \ra{B5a.rrcs.excess.max} points more of the tasks than at the index's own level, and through
duplicate writes that the validation read should have prevented
(Figure~\ref{fig:app-broken}(a)).

\paragraph{Fault injection}
Into the logs of correct runs we planted four kinds of faults and
audited again: a commit that used the answer of its first read although
validation had rejected it (\ri{B5.skip.detected} of
\ri{B5.skip.planted} detected as violations of the safety bound; each
of the others still satisfies the bound, because the fresh answer at
commit is itself approximate), a
commit position moved back in the log (\ri{B5.reorder.detected} of
\ri{B5.reorder.planted}; a move that does not cross a relevant write
is not an error), an answer that contains an item absent from the
commit state (\ri{B5.ghost.detected} of \ri{B5.ghost.planted}), and
distances misreported by the store (\ri{B5.distance.detected} of
\ri{B5.distance.planted}; the auditor recomputes all distances from the
vectors). The unmodified logs of the correct runs have no finding.

%% ====================================================================
\section{Reproducibility}
\label{app:repro}

All experiments, audits, analyses, figures, and tables are produced by
one framework from one command per experiment, with some seed counts
and the MS MARCO subset set on the command line
(\texttt{experiments/vldb/run\_perf.sh}); two steps are manual
(building the patched pgvector, and the dense counterfactual replay of
Figure~\ref{M-fig:pareto} of the paper); the repository is linked in the artifact block of the
paper. Every experiment is a
configuration file that states its question, its grid, and the audits
that run on its logs; every measured number in the paper and in this
appendix is read from the result files by a script (only software
versions, the model names, the machines, and a few counts (data sets,
baselines, models, each checked against the runs) are entered by hand), and the submission build
fails if a number, a figure, or a table is missing. Each run logs, per task, the answers it
reasoned about, the fresh answers, the commit-log positions, and
the decision, and a commit log of all writes; the auditor reconstructs
the exact state at any log position from the workload and the log, so
audits need neither the store nor the index. The PostgreSQL tests
reproduce the findings of \S3 on a stock installation, including
the debugger-based reproduction, and the patch is included. The
reported runs use PostgreSQL \res[00.00]{env.pg} with pgvector
\res[0.0.0]{env.pgvector} and the embedding model \res[a sentence embedding model]{env.embedder}.
Every experiment whose results include a time or a rate (throughput,
latency, critical section, lock wait and occupancy, insertion time)
and both language-model experiments ran on
\res[two servers with NN cores and NNN GB of memory each]{env.machine}: one runs the timing
experiments, store and clients together, and the other the model
and count experiments, whose open model is served by vLLM on \res[a GPU]{env.gpu}.
Every time and every rate in the paper and in this report therefore
comes from these servers. The models that decide are
\res[two language models]{env.llm}; GPT-4o mini answers and judges the downstream
questions on the final memory of every seed. The remaining
experiments only count events (the races of \S3.1, the safety and separation
audits, the conflict-law runs, the counterfactual replays, the index
and tolerance sweeps, and the robustness runs over embedding models,
seeds, retrieval depths, and metrics) and ran on the second of these servers; a third, \res[a server with N CPU cores]{env.vm},
repeats the cost experiment on every workload but MS MARCO. On the cells both ran, the extra calls agree within
\ra{X.machines.offset.bound}\% on average (both ends of the 90\% interval of the mean offset, a random-effects
estimate over \ri{X.machines.cells} cells) and within \ri{X.machines.extra.maxdev}\% on every cell with five or more extra calls per
100 tasks. Numbers are means over seeds; an interval is the 95\% Student $t$ interval of the mean over seeds unless noted. In the generated tables, ``$a\pm h$'' gives the mean $a$ and the half-width $h$ of
that interval (half the range with two seeds), and a mean printed without $\pm$ has an interval that rounds to zero. Every configuration is run with three seeds, except the
model experiments, the conflict-law runs on real streams, the planted-law
simulations, and the exchangeable controls (five seeds), the seed-variance experiment and the lock-scope ablation (ten), the anomaly taxonomy (forty), and the long run and the
Euclidean-distance experiment (one seed each); the cost experiment has eight seeds on L2, L4, and A64 and three elsewhere.

\FloatBarrier
%% ====================================================================


\section{Related Work}
\label{app:related}

\paragraph{Isolation and concurrency control}
Our definition builds on predicate- and state-based accounts of
isolation \cite{DBLP:journals/cacm/EswaranGLT76,DBLP:conf/ds/GrayLPT76,DBLP:journals/jacm/Papadimitriou79,DBLP:books/aw/BernsteinHG87,DBLP:conf/sigmod/BerensonBGMOO95,DBLP:conf/icde/AdyaLO00,Adya1999Weak,DBLP:conf/podc/CrooksPAC17,DBLP:conf/concur/CeroneBG15} and on the
analysis of snapshot isolation \cite{DBLP:journals/tods/FeketeLOOS05,DBLP:conf/sigmod/CahillRF08,DBLP:journals/pvldb/PortsG12}. Phantom protection for ordered
indexes \cite{DBLP:conf/vldb/Mohan90,DBLP:conf/sigmod/MohanL92,DBLP:conf/vldb/Lomet93,DBLP:journals/tods/Graefe10,DBLP:journals/tods/LehmanY81} and main-memory engines \cite{DBLP:conf/sigmod/JordanBB81,DBLP:conf/sosp/TuZKLM13,DBLP:conf/sigmod/NeumannMK15,DBLP:journals/pvldb/LarsonBDFPZ11,DBLP:conf/sigmod/DiaconuFILMSVZ13}, evaluations of
concurrency control \cite{DBLP:journals/pvldb/YuBPDS14,DBLP:journals/pvldb/WuALXP17}, weaker models
\cite{DBLP:journals/pvldb/BailisDFGHS13,DBLP:journals/pvldb/BailisFFGHS14}, and linearizability \cite{DBLP:journals/toplas/HerlihyW90}
assume exact reads. History checkers \cite{DBLP:journals/pvldb/KingsburyA20,DBLP:conf/osdi/TanZMW20,DBLP:journals/pacmpl/BiswasE19} share the intent of our auditor,
which must first decide which answers are legal for an approximate read.
Epsilon serializability \cite{DBLP:conf/sigmod/PuL91,DBLP:journals/tkde/RamamrithamP95}, conits \cite{DBLP:journals/tocs/YuV02,DBLP:journals/pvldb/BailisVFHS12}, relaxed
specifications of concurrent objects \cite{DBLP:conf/opodis/AfekKY10,DBLP:conf/popl/HenzingerKPSS13,DBLP:journals/corr/abs-1804-01018}, and intermediate value
linearizability \cite{DBLP:conf/ppopp/RinbergSBHKRS20,DBLP:conf/disc/RinbergK20,DBLP:journals/pvldb/WangTKZZX23} bound numeric error; admissibility bounds a
ranked set, by a ratio at every rank, in a serial order of
multi-operation transactions.

\paragraph{Multidimensional and nearest-neighbor queries}
R-trees and GiST \cite{DBLP:conf/sigmod/Guttman84,DBLP:conf/vldb/HellersteinNP95,DBLP:conf/vldb/KornackerB95} protect region predicates against phantoms
through predicate and node locks \cite{DBLP:conf/sigmod/KornackerMH97}, granular locks
\cite{DBLP:conf/icde/ChakrabartiM98,DBLP:conf/sigmod/ChakrabartiM99}, and grid cells \cite{DBLP:conf/dasfaa/SongLKY05,DBLP:journals/ijoc/LeeSY07}, for predicates known before the
search over low-dimensional data; linearizable nearest-neighbor search
exists for exact, low-dimensional trees \cite{Chatterjee2015Concurrent,DBLP:conf/spaa/ManoharWB25,Manohar2025New}. Deciding which
open reads an insertion affects is a reverse nearest-neighbor problem
\cite{DBLP:conf/sigmod/KornM00,DBLP:conf/vldb/TaoPL04,DBLP:conf/cikm/SinghFT03,DBLP:journals/corr/abs-1011-4955,DBLP:journals/pvldb/CasanovaEHKNSZ17,DBLP:journals/corr/abs-2606-03225}, akin to continuous query monitoring
\cite{DBLP:conf/sigmod/MouratidisHP05,DBLP:conf/icde/YuPK05,DBLP:conf/icde/XiongMA05,DBLP:conf/icde/BohmOPY07,DBLP:conf/cidr/LuFZSSC25}. Nearest-neighbor and top-$k$ processing
\cite{DBLP:conf/sigmod/RoussopoulosKV95,DBLP:journals/tods/HjaltasonS99,DBLP:conf/stoc/IndykM98,DBLP:journals/csur/IlyasBS08,DBLP:conf/pods/FaginLN01,DBLP:conf/icdt/BeyerGRS99}, duplicate detection
\cite{DBLP:journals/tkde/ElmagarmidIV07,DBLP:conf/icde/ChaudhuriGK06}, and metric integrity constraints
\cite{DBLP:conf/icdt/KaminskyKLNW25} do not treat isolation.

\paragraph{Approximate indexes and vector databases}
Graph and quantization indexes \cite{DBLP:journals/pami/MalkovY20,DBLP:journals/is/MalkovPLK14,DBLP:conf/nips/SubramanyaDSKK19,DBLP:journals/pvldb/FuXWC19,DBLP:journals/pami/JegouDS11,DBLP:journals/tbd/JohnsonDJ21,DBLP:journals/tbd/DouzeGDJSMLHJ26,DBLP:conf/cvpr/BabenkoL12,DBLP:conf/nips/ChenZWLLLYW21,DBLP:journals/pacmmod/GaoL24,DBLP:journals/pacmmod/GaoGXYLW25,DBLP:journals/pvldb/WangXYW21,DBLP:journals/tkde/LiZSWLZL20}, their analyses
\cite{DBLP:conf/nips/DiwanGMMS24,DBLP:conf/nips/IndykX23,DBLP:journals/pvldb/WangWCWPW24,DBLP:conf/sigmod/LiZAH20,DBLP:journals/pacmmod/ChatzakisPP25,DBLP:journals/pvldb/HorchidanZBC25}, filtered, parallel, and GPU variants
\cite{DBLP:journals/pacmmod/PatelKGZ24,DBLP:conf/www/GollapudiKSKBRLZMSSS23,DBLP:conf/ppopp/ManoharSBDGSS24,DBLP:conf/icde/OotomoNNWFW24}, and work that keeps an index fresh under updates
\cite{DBLP:journals/corr/abs-2105-09613,DBLP:conf/sosp/XuLLXCZLYYYCY23,DBLP:journals/corr/abs-2502-13826,DBLP:conf/fast/GuoL26,DBLP:journals/pvldb/YuLGXLZSZLY25,DBLP:journals/corr/abs-2411-00970,DBLP:conf/osdi/MohoneySTCPIRV25,DBLP:journals/corr/abs-2407-07871,DBLP:journals/corr/abs-2507-19802} or versioned
\cite{DBLP:conf/icde/WangHTZZ25,DBLP:journals/corr/abs-2608-27663} reason about the index, not about what a read in a
transaction may return. Vector database systems
\cite{DBLP:conf/sigmod/WangYGJXLWGLXYYZLCLZMGJWX21,DBLP:journals/pvldb/GuoLXYYLCXLLCQWTX22,DBLP:conf/osdi/ZhangXCSXCCHYYYZ23,DBLP:journals/pvldb/WeiWWLZLC20,DBLP:conf/sigmod/YangLFW20,DBLP:journals/pvldb/ChenJZPWWHSWW24,DBLP:conf/sigmod/LiuZCARCXWW25,DBLP:journals/pvldb/SunLPWXLN25,DBLP:journals/pvldb/UpretiSSSBPACSYHDPHMMXDGKGSVBMCCPMKKDSBFV25,DBLP:conf/sigmod/PoundCBGPC25,DBLP:conf/icde/NiuTPC25,DBLP:journals/pvldb/SehgalS25,DBLP:journals/corr/abs-2505-12524,DBLP:conf/icde/ZhangLW24,DBLP:conf/cidr/LiuZJGLW26,DBLP:journals/corr/abs-2608-15994,DBLP:journals/pvldb/OhKKLLZCW26}, surveys
\cite{DBLP:journals/vldb/PanWL24,DBLP:conf/sigmod/PanWL24,DBLP:journals/cogsr/Taipalus24,DBLP:journals/corr/abs-2402-01763,DBLP:journals/corr/abs-2310-11703,DBLP:journals/pvldb/ChronisCPOA25,DBLP:journals/pvldb/WangHLPSX24}, and store documentation
\cite{milvusconsistency,qdrantconsistency,weaviateconsistency,qdrant,weaviate,chroma,lancedb,hnswlib} discuss replica consistency and snapshot
visibility; reliability studies \cite{DBLP:journals/corr/abs-2502-20812,DBLP:journals/corr/abs-2506-02617} call for testing,
public safety analyses cover no vector store yet \cite{jepsen}, and a
practitioner note suspects a repeatable-read problem in pgvector
\cite{songblog}.

\paragraph{Agent memory and transactional agents}
Agent memory architectures
\cite{DBLP:journals/corr/abs-2310-08560,DBLP:conf/uist/ParkOCMLB23,DBLP:conf/ecai/ChhikaraKASY25,DBLP:conf/nips/XuLMGTZ25,DBLP:journals/corr/abs-2501-13956,DBLP:conf/aaai/ZhongGGYW24,DBLP:journals/corr/abs-2507-07957,DBLP:conf/nips/ZhangFWWYY25,DBLP:journals/corr/abs-2505-18279,DBLP:journals/corr/abs-2512-13564,DBLP:journals/tois/ZhangDBMLCZDW25,DBLP:journals/tmlr/SumersYNG24,DBLP:journals/corr/abs-2606-06240},
multi-agent frameworks \cite{DBLP:conf/colm/WuBZWLZJZZLAWBW24,DBLP:conf/iclr/HongZCZCWZWYLZRXWS24,DBLP:conf/nips/LiHIKG23}, and semantic caches
\cite{DBLP:conf/nlposs/Bang23,DBLP:conf/iclr/SchroederDCCLZKKZG26} produce our workload. Analyses of multi-agent failures and data
systems for agents \cite{DBLP:conf/nips/CemriPYACTKPKRZGS25,DBLP:conf/cidr/LiuPSZZACSYSZCCGP26,DBLP:journals/corr/abs-2606-24535,DBLP:journals/corr/abs-2603-10062,DBLP:journals/corr/abs-2608-18092} motivate
transactional agents, and recent proposals bring concurrency control
to tool calls, objects, and workflow resources
\cite{DBLP:journals/corr/abs-2605-17076,DBLP:journals/corr/abs-2606-15376,DBLP:journals/corr/abs-2608-05412,DBLP:journals/corr/abs-2607-23929,DBLP:journals/corr/abs-2608-13900}; none, including the remark that agent searches
resemble predicate reads \cite{DBLP:journals/corr/abs-2610-03140}, defines what a similarity read may
return. We use standard, publicly available embedding models and benchmarks
\cite{DBLP:conf/emnlp/ReimersG19,DBLP:conf/acl/ChenXZLLL24,DBLP:journals/corr/abs-1611-09268,DBLP:conf/iclr/WuWYZCY25,DBLP:journals/is/AumullerBF20,DBLP:conf/nips/SimhadriADBILWLMKDCCMZCSZCYH25,DBLP:conf/cloud/CooperSTRS10,DBLP:conf/acl/MaharanaLTBBF24,DBLP:conf/nips/ZhangZL15,DBLP:conf/iclr/WangSMHLB19}, among them LoCoMo and AG News.

%% ====================================================================
\FloatBarrier
\raggedbottom
\section{Notation and Terms}
\label{app:notation}

This appendix collects every symbol and term of the paper and of this
report, with the place where it is defined. Symbols used only inside
one proof are defined there and not repeated here.

\paragraph{Stores, reads, and answers (\S2)}
\begin{center}\small\setlength{\tabcolsep}{3pt}\renewcommand{\arraystretch}{1.08}
\begin{tabular}{@{}>{\raggedright\arraybackslash}p{0.27\columnwidth}>{\raggedright\arraybackslash}p{0.70\columnwidth}@{}}
\toprule
$S$, $S_0$, $N$ & state (the set of stored items), initial state, number of stored items.\\
$q$, $k$, $\Read(q,k)$ & query vector, number of neighbors, similarity read returning $k$ distinct items in non-decreasing distance from $q$.\\
$\dist$, $\delta_i(S,q)$ & metric on the vectors (chord distance in the evaluation, $\sqrt2$ times the square root of the cosine distance); $i$-th smallest distance from $q$ to $S$.\\
$M$, $\mathit{ef}$ & HNSW degree; search budget, the number of candidates a search keeps.\\
$c$, $\Adm_c(S,q,k)$ & tolerance ($c\ge1$); the set of $c$-admissible answers, lists whose $i$-th item is at most $c$ times as far from $q$ as the true $i$-th neighbor, for every $i$ (Definition~2.1).\\
$\rhoa(A\mid S,q)$ & approximation ratio of the list $A$ in $S$, the least $c$ for which $A$ is $c$-admissible.\\
$\rho$ & a bound on the approximation ratios of the answers a store returns (\S4.2); $\rho=1$ for an exact search. $\rho_{99}$ is the 99th percentile of a ratio over transactions (Table~\ref{M-tab:cost}).\\
$c$-store & a store whose read may return any element of $\Adm_c$ (\S\ref{app:oper}).\\
\bottomrule
\end{tabular}
\end{center}

\paragraph{Transactions and isolation (\S\S2 and~4)}
\begin{center}\small\setlength{\tabcolsep}{3pt}\renewcommand{\arraystretch}{1.08}
\begin{tabular}{@{}>{\raggedright\arraybackslash}p{0.27\columnwidth}>{\raggedright\arraybackslash}p{0.70\columnwidth}@{}}
\toprule
$T$, write set & transaction; its insertions and deletions, buffered and applied atomically at commit. A replacement is a deletion and an insertion.\\
read-think-write & a transaction whose computation between its reads and its writes is a model call.\\
$\tau$, $\tau$-deduplicating insertion & deduplication threshold (stated in cosine distance); a transaction that issues one read and inserts $q$ only if its nearest item is farther than $\tau$.\\
$s$-separated & a state in which any two items are farther apart than $s$ (Proposition~4.2).\\
held answer $A$ & the answer a transaction acted on.\\
execution, $\pi$, $S^\pi_j$ & an initial state and a set of committed transactions; a total order of them; the state after the write sets of the first $j$ transactions in $\pi$.\\
$c$-serializable & there is an order $\pi$ in which every held answer is $c$-admissible in the state just before its transaction (Definition~4.1).\\
\bottomrule
\end{tabular}
\end{center}

\paragraph{Protocol and system (\S\S4.2 and~4.3)}
\begin{center}\small\setlength{\tabcolsep}{3pt}\renewcommand{\arraystretch}{1.08}
\begin{tabular}{@{}>{\raggedright\arraybackslash}p{0.27\columnwidth}>{\raggedright\arraybackslash}p{0.70\columnwidth}@{}}
\toprule
fresh answer $F$, validation read & the answer of the read repeated at commit, under the commit lock.\\
$\sqsubseteq_c$ & validation test: $|A|=|F|$, every item of $A$ is still stored, and $\dist(q,a_i)\le c\,\dist(q,f_i)$ for every $i$.\\
admissibility validation & Algorithm~1 with the test $\sqsubseteq_c$; with an exact validation read it is \emph{certified}.\\
re-execution validation & validation that commits only if $A=F$ (the predicate \pol{exact}).\\
round & one pass through the loop of Algorithm~1: think, validate, and commit or retry; a transaction aborts after 30 failed rounds.\\
commit lock, critical section & the lock under which validation reads and writes run; the time spent holding it (``c.s.'' in the tables).\\
auditor & the tool that replays the commit log, rebuilds the exact state at every commit, and reports the approximation ratio of every held and fresh answer.\\
\bottomrule
\end{tabular}
\end{center}

\paragraph{Conflicts (\S5)}
\begin{center}\small\setlength{\tabcolsep}{3pt}\renewcommand{\arraystretch}{1.08}
\begin{tabular}{@{}>{\raggedright\arraybackslash}p{0.27\columnwidth}>{\raggedright\arraybackslash}p{0.70\columnwidth}@{}}
\toprule
$r_0,\dots,r_k$ & $r_0=0$ and the distances of the exact answer, $r_1<\dots<r_k$; the read's ball has radius $r_k$.\\
in-ball insertion, $t$, $p$ & an insertion at distance $t<r_k$ from $q$; its rank $p$ in the new list.\\
conflict, $\Rc(q)$ & an insertion that makes the held answer inadmissible; the conflict radius, inside which an insertion conflicts (Corollary~5.2).\\
$y$ & position $\ln(r_k/t)$ of an in-ball insertion.\\
overlap & a write overlaps a read if it commits between the read and its validation; $n$ is the number of background writes that overlap a read.\\
$D$, intrinsic dimension & the exponent of the growth of the number of points with distance around $q$; $\lambda s^{D}$ points lie within distance $s$ (\S\ref{app:law-proof}).\\
mass coordinate $\omega$, $\mu_i$, $B_i$ & $\omega=\lambda s^{D}$; masses of the stored neighbors; ratios $\mu_i/\mu_{i+1}$ of consecutive masses (\S\ref{app:law-proof}).\\
background write & a write unrelated to the read, exchangeable with the stored items.\\
sibling write, $\Ds$ & a write about the same thing as the read (a rendering of the same fact, the next fact of the same conversation); its radial density is proportional to $t^{\Ds-1}$ with sibling dimension $\Ds\le D$.\\
$\alpha$, $x$ & dimension ratio $\Ds/D$; $x=c^{-D}$.\\
$\theta_\alpha(p)$, $\varphi_\alpha(p)$, $P_\alpha(c)$ & probability that an in-ball insertion has rank $p$; that it has rank $p$ and leaves the answer admissible; conflict probability of one in-ball insertion (Theorem~5.3).\\
$\sigma$, two scales & share of sibling writes among in-ball insertions; a mixture of two sibling dimensions with their shares $\sigma_j$, fitted when one dimension does not describe a workload.\\
exchangeable model, closed form, semi-parametric form & the law with every write treated as background; Theorem~5.3 with the pooled intrinsic dimension; the sibling law evaluated on the read's actual distances (\S5.2).\\
\bottomrule
\end{tabular}
\end{center}

\paragraph{The lag kernel (\S5.3 and \S\ref{app:est})}
\begin{center}\small\setlength{\tabcolsep}{3pt}\renewcommand{\arraystretch}{1.08}
\begin{tabular}{@{}>{\raggedright\arraybackslash}p{0.27\columnwidth}>{\raggedright\arraybackslash}p{0.70\columnwidth}@{}}
\toprule
lag $\ell$, lag bucket $b$ & distance in the stream between the reading item and the source of a write, in the same stream (the same conversation or feed); a range of lags.\\
$h_\ell$, $D_\ell$; $\hat h_b$, $\hat D_b$ & probability that the item at lag $\ell$ falls into the read's ball, and its sibling dimension; their estimates per lag bucket.\\
$o_\ell$, $o_b$, $\nu$ & expected number of overlapping writes at lag $\ell$ (or in bucket $b$); expected in-ball insertions per read.\\
pooled parameters & sibling dimensions and shares fitted over all lags of one run.\\
static stream & the estimate of the kernel from the stream alone, without any execution.\\
transfer & a prediction for a run from parameters fitted on another run (another concurrency level, $k$, time, collection, or the static stream).\\
$w_{\mathrm{bg}}$, $f_{\mathrm{bg}}$ & background share of in-ball insertions, fixed from first principles; background density of $y$ (\S\ref{app:est}).\\
\bottomrule
\end{tabular}
\end{center}

\paragraph{Evaluation (\S6)}
\begin{center}\small\setlength{\tabcolsep}{3pt}\renewcommand{\arraystretch}{1.08}
\begin{tabular}{@{}>{\raggedright\arraybackslash}p{0.27\columnwidth}>{\raggedright\arraybackslash}p{0.70\columnwidth}@{}}
\toprule
fan-out workload & 600 facts, each written by one, two, or four agents as perturbed copies (\S3.1).\\
L2, L4, L8 & LoCoMo with teams of 2, 4, or 8 agents per conversation; $g$ is the team size.\\
A16, A64, A64k50, A256, A64-long & AG News with 16, 64, or 256 agents; A64k50 reads $k=50$; A256 writes 12{,}000 items, A64-long 30{,}000.\\
LME4, QQP64, MSM64 & LongMemEval with teams of four agents; Quora questions with 64 agents; MS MARCO passages with 64 agents.\\
extra calls & model calls beyond one per task, per 100 tasks.\\
decision proxies & \pol{dup} (is the nearest item within $\tau$), \pol{tag} (majority label of the $k$ items), \pol{link} (identity of the nearest item), \pol{top3} (identity of the three nearest), \pol{set} (identity of all $k$).\\
predicates & \pol{none} (no isolation), \pol{exact} (re-execution), \pol{set}, \pol{overlap} (at least 90\% common identifiers), \pol{objects} (the items read still exist), the $k$-th distance (the radius), the first distance, and the sum of distances only.\\
index floor & the share of transactions whose decision on the fresh answer differs from that on the exact answer at commit.\\
\bottomrule
\end{tabular}
\end{center}

\FloatBarrier
%% ====================================================================
\bibliographystyle{ACM-Reference-Format}
\bibliography{pvldb20-ariadne}

\end{document}
